\scanpage{196}
\margin{195}
ALGÈBRES de LIE SEMI-SIMPLES (7) \hfill Tribu n° 59

mettre l'exer. 1, (p. 36), en c) de la prop. 11, et l'exer. 2 dans le texte.

Ajouter un utile rabiot sur les algèbres de la forme, (commutative), $\times\underline{\mathrm{sl}}_2$, à savoir les généralisations des prop. 5 et 6.

\S 3 -- \uuline{Le système de racines...}

\uline{n° 1} -- Définir ``déploiement'', (i.e. couple $(\underline g,\underline h)$), ``déployable''. Dire que $(*)$ de la déf. 1 signifie que $\operatorname{ad}(x)$ est diagonalisable, (on sait qu'il est semi-simple). Ajouter ``semi-simple'' à la l.2 bas, (p. 37).

Remplacer ``canonique'' par ``diagonale'', et ``abélienne'' par ``commutative'', ousqui faut.

\uline{n° 2} -- Il s'agit du sous-espace ``primaire'' à la p. 38, l.4. A la l.5 dire : ``comme $\operatorname{ad}(\underline h)$ est semi-simple, $\underline g^\alpha$ est l'ensemble...'', et ne pas répéter ça à la l.16. Dans l'énoncé du th. 2 dire que $\underline h^\alpha$ est un sous-espace de dim. 1 de $\underline h$. Y expliciter l'isom. $\varphi_\alpha$ dans (iii). Ajouter ``et si'' dans le (iv).

\uline{n° 3} -- Après le premier alinéa de la p. 41, démontrer le th. 3, c), (cf. p. 44), ce qui règle son compte au th. 3. Plus bas utiliser l'isom. $\varphi_\alpha$. Commencer la p. 42 par ``Soit...''. A la l.3, (p. 42), qualifier $W(2)$ de ``module simple''. Ce qui est entre parenthèses en-dessous est maintenant fait au \S 2. Dire, (p. 42, l.12), que $R$ est un système de racines `` de $\underline h$ ''.

Le haut de la p. 44 remonte p. 41.

\uline{n° 4} -- Titre : terminologie et inventaire.

\uline{n° 5} -- A la l.2 bas de la p. 45 dire : ``Dans tout couple $(\alpha,-\alpha)$ de racines opposées, choisissons arbitrairement une racine $\alpha$ ; pour chaque racine $\alpha$ ainsi choisie, choisissons... $X_\alpha$...''.
